\begin{titlepage}
  \centering
  \vspace*{18mm}
  {\Large\sffamily\color{CasuloTeal} Guilherme Moura Fernandes\par}
  \vfill
  {\Huge\sffamily\bfseries\color{CasuloNavy}
    Princípio da Interlocução\par}
  \vspace{8mm}
  {\LARGE\sffamily
    a Linguagem T como arquiteta, auditora e juíza da integridade semântica em
    modelos de treinamento de redes neurais artificiais e a decorrente
    compressão física da informação preservada\par}
  \vfill
  {\large\sffamily 2026\par}
\end{titlepage}

\cleardoublepage
\chapter*{Resumo}
\addcontentsline{toc}{chapter}{Resumo}
Este exemplar editorial demonstra a estrutura tipográfica destinada à análise
do Princípio da Interlocução. O manuscrito final deverá reconstruir, a partir de
artefatos hashados, a autoridade causal da Linguagem T, a equivalência
operacional e a contabilidade multidimensional do trabalho, preservando limites
e resultados negativos. Nenhum valor deste parágrafo substitui a análise dos
dados primários.

\noindent\textbf{Palavras-chave:} Linguagem T; integridade semântica; autoridade
causal; treinamento neural; compiladores; reprodutibilidade.

\selectlanguage{english}
\chapter*{Abstract}
\addcontentsline{toc}{chapter}{Abstract}
This editorial sample establishes the typographic structure for an evidence-led
analysis of the Interlocution Principle. The final manuscript must reconstruct
T Language causal authority, operational equivalence, and multidimensional work
accounting from hashed artifacts while preserving limitations and negative
results.

\noindent\textbf{Keywords:} T Language; semantic integrity; causal authority;
neural training; compilers; reproducibility.
\selectlanguage{brazil}

\tableofcontents
\listoffigures
\listoftables
